
\chapter{Richard Schoen (2017)}
\index{Schoen, Richard}

\begin{quote}
\it ``Richard Schoen's work has transformed our understanding of the relationship between geometry and analysis. His proofs of the positive mass theorem and the Yamabe problem, together with his fundamental contributions to minimal surface theory and harmonic maps, have established him as one of the leading geometers of our time.'' --- Wolf Prize Citation, 2017
\end{quote}

\section{Early Life and Career}

Richard Michael Schoen was born on October 23, 1950, in Fort Wayne, Indiana. He studied at the University of Michigan and then at Stanford University, where he received his Ph.D.\ in 1977 under the supervision of Robert Gunning. He held positions at the University of California, Berkeley, Stanford University, and the University of California, Irvine. He has been a professor at Stanford since 1986. He was awarded the Wolf Prize in Mathematics in 2017.

\section{Main Contributions}

Richard Schoen (born 1950) is an American mathematician whose work in geometric analysis has produced some of the most important results in modern geometry. He was awarded the Wolf Prize in Mathematics in 2017 alongside Charles Fefferman.

The Wolf Prize citation reads: ``for their contributions to analysis and geometry.''

Schoen's core contributions include:
\begin{itemize}
    \item \textbf{Positive Mass Theorem:\index{positive mass theorem}} A proof (with Yau) that the ADM mass of an asymptotically flat manifold with nonnegative scalar curvature is nonnegative, with equality only for Euclidean space.
    \item \textbf{The Yamabe Problem:} The resolution of the Yamabe problem: every compact Riemannian manifold of dimension $n \geq 3$ admits a metric of constant scalar curvature in its conformal class.
    \item \textbf{Minimal Surface Theory:} Foundational results on the existence, regularity, and stability of minimal surfaces, including the proof (with Yau) that any compact 3-manifold with positive scalar curvature contains a stable minimal surface.
    \item \textbf{Harmonic Maps:} The Schoen--Uhlenbeck theory of harmonic maps, including regularity and compactness results.
    \item \textbf{Conformal Geometry:} The Schoen--Simon--Yau estimates for minimal surfaces and the classification of stable minimal hypersurfaces.
\end{itemize}

\section{Origin of the Problem}

\subsection{The Positive Mass Problem}

In general relativity, the total mass (energy) of an isolated gravitational system is defined by the ADM mass. A fundamental physical question: can this mass be negative? Negative mass would allow instability and physically pathological behavior. Proving positivity was essential for the stability of Minkowski spacetime.

The problem reduces to: on an asymptotically flat Riemannian 3-manifold with nonnegative scalar curvature, prove the ADM mass is nonnegative.

\subsection{The Yamabe Problem}

Given a compact Riemannian manifold $(M, g)$ of dimension $n \geq 3$, the Yamabe problem asks: does there exist a conformal metric $\tilde g = u^{4/(n-2)} g$ with constant scalar curvature?

This is the Yamabe equation:
\[
-\frac{4(n-1)}{n-2} \Delta u + \Scal_g u = \lambda u^{(n+2)/(n-2)}
\]

Trudinger solved the case when the Yamabe invariant is nonpositive. Aubin solved the case when $n \geq 6$ and $M$ is not locally conformally flat. The remaining difficult case (locally conformally flat or $n = 3,4,5$) required Schoen's deep analytic and geometric insights.

\subsection{Minimal Surfaces in 3-Manifolds}

A minimal surface is a surface that locally minimizes area. In a 3-manifold with positive scalar curvature, stable minimal surfaces have restricted topology. Schoen and Yau used this to prove topological restrictions on manifolds with positive scalar curvature.

\section{How It Was Solved and the Intuitions/Tools Needed}

\subsection{The Proof of the Positive Mass Theorem}

Schoen and Yau's proof uses minimal surface techniques:
\begin{enumerate}
    \item Suppose the ADM mass $m < 0$. One can deform the metric in a compact region to produce a metric with positive scalar curvature outside but zero scalar curvature inside, with $m$ still negative.
    \item Using the negative mass, construct a non-compact stable minimal surface.
    \item Show that such a surface must have positive Gauss curvature, contradicting the Gauss--Bonnet theorem (the surface would have the topology of a sphere but the metric would be complete and noncompact).
\end{enumerate}

\subsection{The Proof of the Yamabe Problem}

Schoen's proof uses the following strategy:
\begin{enumerate}
    \item If the Yamabe invariant $Y(M)$ is less than the Yamabe invariant $Y(S^n)$ of the sphere, then Aubin's methods apply and the problem is solved.
    \item If $Y(M) = Y(S^n)$, one must show that equality cannot occur unless $(M,g)$ is conformally equivalent to the sphere.
    \item Schoen developed the positive mass theorem as a tool: he showed that if $Y(M) = Y(S^n)$, then a certain test function can be constructed whose conformal Laplacian has positive mass, contradicting the positive mass theorem.
\end{enumerate}

This beautiful connection between the Yamabe problem (a geometric PDE) and the positive mass theorem (a general relativity result) was a revelation.

\section{Mathematical Theory}

\subsection{Positive Mass Theorem}

\begin{definition}[ADM Mass]
Let $(M, g)$ be an asymptotically flat Riemannian $n$-manifold (with $n \geq 3$). In asymptotic coordinates $x$, the metric satisfies $g_{ij} = \delta_{ij} + O(|x|^{1-n/2})$. The \textbf{ADM mass} is:
\[
m = \frac{1}{2(n-1)\omega_{n-1}} \lim_{r\to\infty} \int_{S_r} \sum_{i,j} (\partial_j g_{ij} - \partial_i g_{jj}) \frac{x^i}{r} \, dS
\]
\end{definition}

\begin{theorem}[Schoen--Yau, 1979]
Let $(M,g)$ be an asymptotically flat Riemannian $3$-manifold with nonnegative scalar curvature. Then $m \geq 0$, with equality if and only if $(M,g)$ is isometric to Euclidean space $(\R^3, \delta)$.
\end{theorem}

\begin{corollary}
In general relativity, the total energy of an isolated gravitational system is nonnegative, confirming a fundamental physical stability condition.
\end{corollary}

\subsection{The Yamabe Problem}

\begin{definition}[Yamabe Invariant]
For a compact Riemannian manifold $(M,g)$ of dimension $n \geq 3$, the \textbf{Yamabe invariant} is:
\[
Y(M,g) = \inf_{u > 0} \frac{\int_M (c_n |\nabla u|^2 + \Scal_g u^2) \, dV}{\left(\int_M u^{2n/(n-2)} \, dV\right)^{(n-2)/n}}
\]
where $c_n = 4(n-1)/(n-2)$.
\end{definition}

\begin{theorem}[Schoen, 1984]
Every compact Riemannian manifold of dimension $n \geq 3$ admits a conformal metric of constant scalar curvature.
\end{theorem}

The key cases solved by Schoen:
\begin{itemize}
    \item \textbf{Case $Y(M,g) \leq 0$:} Solved by Trudinger using subcritical approximation.
    \item \textbf{Case $Y(M,g) > 0$ and $n \geq 6$, $M$ not locally conformally flat:} Solved by Aubin using test functions.
    \item \textbf{Case $Y(M,g) > 0$ and $n = 3,4,5$ or $M$ locally conformally flat:} Solved by Schoen using the positive mass theorem.
\end{itemize}

\subsection{Topology of Manifolds with Positive Scalar Curvature}

\begin{theorem}[Schoen--Yau, 1979]
Let $M$ be a compact $3$-manifold with positive scalar curvature. Then $\pi_1(M)$ contains no subgroup isomorphic to the fundamental group of a compact surface of non-negative Euler characteristic.
\end{theorem}

This result was proved by constructing a stable minimal surface and analyzing its geometry. It was later extended to higher dimensions (Schoen--Yau, Gromov--Lawson) and is central to the theory of positive scalar curvature.

\subsection{Harmonic Maps}

\begin{definition}[Harmonic Map]
A map $f: M \to N$ between Riemannian manifolds is \textbf{harmonic} if it is a critical point of the energy functional:
\[
E(f) = \frac{1}{2} \int_M |df|^2 \, dV
\]
The Euler--Lagrange equation is $\tau(f) = 0$, where $\tau(f)$ is the tension field.
\end{definition}

\begin{theorem}[Schoen--Uhlenbeck, 1982]
Let $M$ be a compact Riemannian $m$-manifold and $N$ a compact Riemannian manifold. A weakly harmonic map $f: M \to N$ is smooth outside a closed set $\Sigma$ of Hausdorff codimension at least $3$. In particular, for $m = 2$, all weakly harmonic maps are smooth.
\end{theorem}

\subsection{The Positive Mass Theorem: Full Statement}

\begin{theorem}[Positive Mass Theorem; Schoen--Yau, 1979 (Spin Case); Witten, 1981 (General Case)]
Let $(M, g)$ be an asymptotically flat Riemannian $n$-manifold ($n \geq 3$) with nonnegative scalar curvature $\Scal_g \geq 0$. If $M$ is spin, then the ADM mass satisfies:
\[
m_{\text{ADM}} \geq 0
\]
with equality if and only if $(M, g)$ is isometric to $(\R^n, \delta)$. In the general (non-spin) case, the theorem holds with the additional hypothesis that the manifold is orientable and the scalar curvature satisfies $\Scal_g \geq 0$.
\end{theorem}

\begin{corollary}
The ADM mass of an asymptotically flat $3$-manifold with $\Scal_g \geq 0$ and $M$ spin satisfies:
\[
m_{\text{ADM}} = \frac{1}{16\pi} \lim_{r \to \infty} \int_{S_r} \sum_{i,j} (\partial_j g_{ij} - \partial_i g_{jj}) \frac{x^i}{r} \, dS \geq 0
\]
\end{corollary}

\begin{remark}
Witten's proof (1981) uses spinors and the Dirac operator. The key identity relates the mass to an integral of $|\nabla \psi|^2 + \frac{1}{4}\Scal |\psi|^2$ for a harmonic spinor $\psi$, which is nonnegative when $\Scal \geq 0$. Schoen and Yau's original proof used minimal surfaces and is valid without the spin assumption in dimension $3$.
\end{remark}

\subsection{Schoen--Yau Iteration Technique}

\begin{definition}[Prescribed Scalar Curvature Equation]
Given a compact Riemannian manifold $(M, g)$ of dimension $n \geq 3$ and a smooth function $f$ on $M$, the \textbf{prescribed scalar curvature equation} is the semilinear elliptic PDE:
\[
-\frac{4(n-1)}{n-2} \Delta_g u + \Scal_g \, u = f \cdot u^{(n+2)/(n-2)}
\]
A solution $u > 0$ gives the conformal metric $\tilde{g} = u^{4/(n-2)} g$ with $\Scal_{\tilde{g}} = f$.
\end{definition}

\begin{theorem}[Schoen--Yau Iteration, 1988]
The \textbf{Schoen--Yau iteration} for the Yamabe problem proceeds as follows:
\begin{enumerate}
    \item Start with a conformal factor $u_0 > 0$ and define $g_0 = u_0^{4/(n-2)} g$.
    \item At each step, solve the linearized equation:
    \[
    L_{g_k} v_k = -c_n \Delta_{g_k} v_k + \Scal_{g_k} v_k = \lambda_k v_k^{(n+2)/(n-2)}
    \]
    where $L_{g_k}$ is the conformal Laplacian of $g_k$.
    \item Use the Yamabe functional as a Lyapunov function:
    \[
    Q(g_k) = \frac{\int_M \Scal_{g_k} \, dV_{g_k}}{\left(\int_M dV_{g_k}\right)^{(n-2)/n}}
    \]
    which decreases under the iteration: $Q(g_{k+1}) \leq Q(g_k)$.
    \item If the Yamabe invariant $Y(M, g) > 0$ and equals $Y(S^n) = n(n-1)\omega_n^{2/n}$, construct test functions $u_\epsilon$ concentrating at a point and show that the limiting mass is positive, contradicting the positive mass theorem unless $(M, g)$ is conformally equivalent to $S^n$.
\end{enumerate}
\end{theorem}

\begin{remark}
The iteration scheme is a gradient flow for the Yamabe functional. The key analytic input is the Moser--Trudinger inequality (in the locally conformally flat case) and the positive mass theorem (for the borderline case $Y(M, g) = Y(S^n)$). The iteration converges because the Yamabe functional is bounded below and the conformal Laplacian satisfies a maximum principle.
\end{remark}

\subsection{Minimal Surface Methods}

\begin{definition}[Stable Minimal Surface]
A two-sided minimal surface $\Sigma$ in a Riemannian $3$-manifold $(M, g)$ is \textbf{stable} if the second variation of area is nonnegative for all compactly supported variations of $\Sigma$:
\[
\delta^2 A(\Sigma)(\phi, \phi) = \int_\Sigma \bigl(|\nabla \phi|^2 - (\text{Ric}_M(\nu, \nu) + |A|^2) \phi^2\bigr) \, d\Sigma \geq 0
\]
for all $\phi \in C_c^\infty(\Sigma)$, where $\nu$ is the unit normal to $\Sigma$ and $|A|^2$ is the squared norm of the second fundamental form.
\end{definition}

\begin{theorem}[Schoen--Yau, 1979; Minimal Surface Methods for Positive Mass]
The minimal surface method for the positive mass theorem proceeds as follows:
\begin{enumerate}
    \item Suppose $(M, g)$ is an asymptotically flat $3$-manifold with $\Scal_g \geq 0$ and $m_{\text{ADM}} < 0$.
    \item By the assumption $m_{\text{ADM}} < 0$, one can construct a complete, noncompact stable minimal surface $\Sigma$ in $(M, g)$.
    \item The stability inequality, combined with the Gauss equation, gives:
    \[
    \int_\Sigma \bigl(|\nabla \phi|^2 - (\text{Ric}_M(\nu, \nu) + |A|^2) \phi^2\bigr) \, d\Sigma \geq 0
    \]
    \item Using $\Scal_g = 2K_\Sigma - |A|^2 + H^2 - \text{Ric}_M(\nu,\nu)$ (where $K_\Sigma$ is the Gauss curvature of $\Sigma$ and $H = 0$ for minimal surfaces), and $\Scal_g \geq 0$, one obtains:
    \[
    \int_\Sigma |A|^2 \phi^2 \, d\Sigma \leq \int_\Sigma \bigl(|\nabla \phi|^2 + 2K_\Sigma \phi^2\bigr) \, d\Sigma
    \]
    \item Choosing $\phi$ to be a cutoff function and using the fact that $\Sigma$ is complete and noncompact, one derives that $K_\Sigma \leq 0$ outside a compact set.
    \item By the Gauss--Bonnet theorem, $\Sigma$ must have the topology of a plane, and the curvature estimate leads to a contradiction with the completeness of $\Sigma$.
\end{enumerate}
\end{theorem}

\begin{remark}
The minimal surface method was later extended by Schoen and Yau to prove topological restrictions on $3$-manifolds with positive scalar curvature. Any incompressible surface in such a manifold must have non-positive Euler characteristic. This was further generalized to higher dimensions by Gromov and Lawson.
\end{remark}

\begin{theorem}[Schoen--Yau, 1979; Topological Restriction via Minimal Surfaces]
Let $M$ be a compact orientable $3$-manifold with $\Scal_g > 0$. Then:
\begin{enumerate}
    \item $M$ contains no immersed incompressible torus or Klein bottle.
    \item Any embedded incompressible surface $\Sigma \subset M$ must satisfy $\chi(\Sigma) < 0$ (i.e., $\Sigma$ is a surface of genus $\geq 2$).
    \item The fundamental group $\pi_1(M)$ contains no free abelian subgroup of rank $2$.
\end{enumerate}
\end{theorem}

\section{Impact on Science and Technology}

\subsection{In Mathematics}

\begin{enumerate}
    \item \textbf{Geometric Analysis:} Schoen's proofs demonstrated the power of combining PDE techniques with differential geometry.
    \item \textbf{General Relativity:} The positive mass theorem is a cornerstone of mathematical relativity.
    \item \textbf{Minimal Surface Theory:} Schoen's work with Yau established the modern approach to minimal surfaces in 3-manifolds.
    \item \textbf{Conformal Geometry:} The solution of the Yamabe problem completed the program begun by Yamabe, Trudinger, and Aubin.
    \item \textbf{Harmonic Map Theory:} The Schoen--Uhlenbeck theory is the standard regularity theory for harmonic maps.
\end{enumerate}

\subsection{In Physics}

\begin{enumerate}
    \item \textbf{General Relativity:} The positive mass theorem is essential for the stability of Minkowski space and black hole thermodynamics.
    \item \textbf{String Theory:} Calabi--Yau compactifications rely on many of the geometric analysis techniques Schoen developed.
\end{enumerate}

\section{Frequently Asked Questions}

\subsection*{Q1: What is the positive mass theorem?}
It states that the total gravitational energy of an isolated system (the ADM mass) is always nonnegative if the local energy density (scalar curvature) is nonnegative. This confirms that negative mass cannot exist in classical general relativity.

\subsection*{Q2: What is the Yamabe problem?}
The question: given a compact Riemannian manifold, can we find a conformal metric (one that scales the original metric pointwise) with constant scalar curvature? Schoen proved the answer is always yes.

\subsection*{Q3: How is the Yamabe problem related to the positive mass theorem?}
In the most difficult case, where the Yamabe invariant equals that of the sphere, Schoen showed that if no solution existed, one could construct a metric violating the positive mass theorem. Thus the positive mass theorem forces the existence of a solution.

\subsection*{Q4: What do minimal surfaces have to do with scalar curvature?}
Stable minimal surfaces in a 3-manifold with positive scalar curvature must be spheres or planes. This restricts the topology of the ambient manifold and was used by Schoen and Yau to prove the positive mass theorem and topological restrictions.

\subsection*{Q5: What should an undergraduate study to understand Schoen's work?}
Differential geometry (Riemannian geometry, curvature, geodesics), PDE theory (elliptic equations, Sobolev spaces), and minimal surface theory. Recommended: \textit{Riemannian Geometry} (do Carmo), \textit{Differential Geometry of Curves and Surfaces} (do Carmo), and \textit{Partial Differential Equations} (Evans).

\section{Related Chapters}
See also Chapter~23 (De Giorgi) on regularity theory and Chapter~50 (Yau) on differential geometry and the positive mass theorem.

\section*{Exercises for the Reader}
\addcontentsline{toc}{section}{Exercises}

\begin{enumerate}
    \item [B] Compute the ADM mass for the Schwarzschild metric $g = (1 + 2m/r)^{-1} dr^2 + r^2 d\Omega^2$.
    \item [I] Show that the Yamabe invariant $Y(S^n) = n(n-1)\omega_n^{2/n}$.
    \item [A] Prove that a stable minimal surface in a 3-manifold with positive scalar curvature must be a sphere.
    \item [B] Show that the Yamabe equation is the Euler--Lagrange equation for the Yamabe functional.
    \item [B] Derive the Gauss--Bonnet theorem for a compact surface.
    \item [I] Prove that a harmonic map from a compact surface to a sphere of nonpositive curvature is constant.
    \item [B] Show that the conformal Laplacian transforms conformally.
\end{enumerate}

\section*{Timeline}
\addcontentsline{toc}{section}{Timeline}

\begin{tabularx}{\linewidth}{|p{1.5cm}|>{\raggedright\arraybackslash}X|}
\hline
\textbf{Year} & \textbf{Event} \\ \hline
1950 & Born in Celina, Ohio \\ \hline
1971 & B.S.\ from University of Dayton \\ \hline
1975 & Ph.D.\ from Stanford (advisor: Leon Simon) \\ \hline
1979 | Positive mass theorem (with Yau) \\ \hline
1982 | Regularity of harmonic maps (with Uhlenbeck) \\ \hline
1984 | Solution of the Yamabe problem \\ \hline
1988 | MacArthur Fellowship \\ \hline
1989 | Bocher Prize of the AMS \\ \hline
2005 | Moves to UC Irvine \\ \hline
2007 | Moves to Stanford University \\ \hline
2017 | Wolf Prize in Mathematics \\ \hline
\end{tabularx}

\section*{Glossary}
\addcontentsline{toc}{section}{Glossary}

\begin{description}
    \item[ADM mass] The total energy of an asymptotically flat gravitational system.
    \item[Asymptotically flat] A Riemannian manifold that approaches Euclidean space at infinity.
    \item[Conformal metric] A metric $\tilde g = e^{2f} g$ that scales the original metric pointwise.
    \item[Harmonic map] A map between Riemannian manifolds minimizing the energy functional.
    \item[Minimal surface] A surface with zero mean curvature, locally area-minimizing.
    \item[Scalar curvature] The simplest curvature invariant, obtained by fully contracting the Riemann tensor.
    \item[Yamabe invariant] A conformal invariant measuring the best constant in the Sobolev inequality.
\end{description}

\section{Citations}\subsection*{Primary Sources}
\begin{enumerate}
    \item R.\ Schoen and S.-T.\ Yau, \textit{On the proof of the positive mass conjecture in general relativity}, Comm.\ Math.\ Phys.\ 65 (1979), 45--76.
    \item R.\ Schoen, \textit{Conformal deformation of a Riemannian metric to constant scalar curvature}, J.\ Diff.\ Geom.\ 20 (1984), 479--495.
    \item R.\ Schoen and S.-T.\ Yau, \textit{Existence of incompressible minimal surfaces and the topology of three-manifolds with nonnegative scalar curvature}, Ann.\ Math.\ 110 (1979), 127--142.
    \item R.\ Schoen and K.\ Uhlenbeck, \textit{A regularity theory for harmonic maps}, J.\ Diff.\ Geom.\ 17 (1982), 307--335.
    \item R.\ Schoen and S.-T.\ Yau, \textit{On the structure of manifolds with positive scalar curvature}, Manuscr.\ Math.\ 28 (1979), 159--183.
\end{enumerate}

\subsection*{Expository Works}
\begin{enumerate}
    \item J.\ M.\ Lee and T.\ H.\ Parker, \textit{The Yamabe problem}, Bull.\ AMS 17 (1987), 37--91.
    \item R.\ Schoen, \textit{Variational theory for the total scalar curvature functional for Riemannian metrics and related topics}, Proc.\ ICM 1986.
    \item R.\ Schoen and S.-T.\ Yau, \textit{Lectures on Differential Geometry}, International Press, 1994.
\end{enumerate}

